\chapter{Examples of ABC and Quillen model categories}
\label{chap:examples}

\section{Topological spaces}
\label{sec:exampletopspaces}
\index{$Top$}
In the category of topological spaces $Top$, we denote $B^A$ for the space of continuous maps $A \ra B$, with the compact-open topology. The exponential map $f \mapsto (a \mapsto (b \mapsto f(a, b)))$ defines a bijection
\begin{equation}
\label{eqn:exp}
Top(A \times B, C) \cong Top(A, C^B)
\end{equation}
for any topological spaces $A, B, C$ with $B$ locally compact and Hausdorff separated. 

%In \definitionref{defn:homotopicmaps} further down, for a category ${\cat M}$ with a class of weak equivalences ${\cat W}$ we will define two maps $f, g \colon A \ra B$ to be homotopic $f \simeq g$ if they have the same image in ${\bf ho}{\cat M}$. In the case of $Top$, we will proceed in the reverse way: we will first define $\simeq$ as the classic homotopy relation of maps in $Top$. We will then define the weak equivalences to be the homotopy equivalences, and show in \propositionref{prop:homotopyconsistence} that two maps $f, g \colon A \ra B$ in $Top$ have the same image in ${\bf ho}(Top)$ iff $f \simeq g$. This ensures that the classic homotopy $\simeq$ in $Top$ also satisfies \definitionref{defn:homotopicmaps}.

Denote $I$ the unit interval in $Top$. For any space $A$, we denote $i_0, i_1 \colon A \ra I \times A$ the inclusions $i_k(a) = (k, a)$ for $k = 0, 1$, and denote $p \colon I \times A \ra A$ the second factor projection. We also denote $p_0, p_1 \colon A^I \ra A$ the evaluation maps $p_k(f) = f(k)$, and $i \colon A \ra A^I$ the constant-value map $i(a)(t) = a$. 

\begin{defn}
\label{defn:homotopy}
Suppose that $f_0, f_1 \colon A \ra B$ is a pair of maps in $Top$. A homotopy $h \colon f_0 \simeq f_1$ (also denoted $f_0 \xymatrix{\ar[r]|h &} f_1$) is a map $h \colon I \times A \ra B$ with $hi_k = f_k$. Using the bijection \eqref{eqn:exp}, a homotopy $h \colon I \times A \ra B$ amounts to a map $h^{'} \colon A \ra B^I$ with $p_kh^{'} = f_k$.
\end{defn}

\begin{defn}
\label{defn:homotopyundera}
Suppose that $f_0, f_1$ are two maps under $A$, meaning that $f_kb=c$ for $k = 0, 1$.
\[
\xymatrix{
  & A \ar[dl]_b \ar[dr]^{c} & \\
  B \ar@<.5ex>[rr]^{f_0} \ar@<-.5ex>[rr]_{f_1} & & C
}
\]
A homotopy under $A$, denoted $h \colon f_0 \overset{A}{\simeq} f_1$ is a homotopy $h \colon f_0 \simeq f_1$ which is {\it constant} when restricted to $I \times A$, meaning that $h \comp (I \times b) = cp$, where $p \colon I \times A \ra A$ is the projection.
\end{defn}

\begin{defn}
\label{defn:homotopyovera}
Suppose that $f_0, f_1$ are two maps over $A$, meaning that $cf_k=b$ for $k = 0, 1$.
\[
\xymatrix{
  B \ar@<.5ex>[rr]^{f_0} \ar@<-.5ex>[rr]_{f_1} \ar[dr]_b & & C \ar[dl]^c \\
  & A &
}
\]
A homotopy over $A$, denoted $h \colon f_0 \underset{A}{\simeq} f_1$ is a homotopy $h \colon f_0 \simeq f_1$ which is {\it constant} when corestricted to $C$, meaning that $ch = bp$, where $p \colon I \times B \ra B$ is the projection.
\end{defn}

Homotopy $\simeq$ and its relative counterparts $\overset{A}{\simeq}$, $\underset{A}{\simeq}$ are equivalence relations.

\begin{defn}
\label{defn:homotequiv}
A map $f \colon A \ra B$ is a {\it homotopy equivalence} if it admits a map $g \colon B \ra A$ with $fg \simeq 1_B$, $gf \simeq 1_A$.
\end{defn}

Homotopy equivalences satisfy the 2 out of 3 axiom, and are closed under small sums and small products.

\index{Hurewicz (co)fibration}
\begin{defn}
\label{defn:hep}
\mbox{}
\begin{enumerate}
\item
A map $i \colon A \ra B$ is a {\it Hurewicz cofibration} if it has the left lifting property with respect to all maps of the form $p_0 \colon C^I \ra C$.
\item
A map $p \colon A \ra B$ is a {\it Hurewicz fibration}  if it has the right lifting property with respect to all maps of the form $i_0 \colon C \ra I \times C$.
\end{enumerate}
\end{defn}

By definition, a Hurewicz cofibration (or fibration) is {\it trivial} if it is also a homotopy equivalence. The following result is immediate.


\begin{lem}
\label{lem:hurcomp}
\mbox{}
\begin{enumerate}
\item
(Trivial) Hurewicz cofibrations are closed under compositions and small sums.
\item
(Trivial) Hurewicz fibrations are closed under compositions and small products. $\square$
\end{enumerate}
\end{lem}

Using the bijection \eqref{eqn:exp} we see that $i \colon A \ra B$ is a Hurewicz cofibration iff it has the {\it homotopy extension property (HEP)}, meaning that for any maps $f, h$ with $fi = hi_0$ 
\begin{equation}
\label{eqn:hep}
\xymatrix{
  A \ar[r]^-{i_0} \ar[d]_i & I \times A \ar[d]_{I \times i} \ar[rdd]^h & \\
  B \ar[r]^-{i_0} \ar[drr]_f & I \times B \ar@{-->}[rd]^(.3){h^{'}} & \\
  && C
}
\end{equation}
there exists a homotopy $h^{'}$ keeping the diagram commutative. Dually, a map $p \colon A \ra B$ is a Hurewicz fibration iff it has the {\it homotopy lifting property (HLP)}, meaning that for any maps $f, h$ with $pf = p_0h$ 
\begin{equation}
\label{eqn:hlp}
\xymatrix{
  C \ar[rrd]^f \ar[rdd]_h \ar@{-->}[rd]^(.9){h^{'}} && \\
  & A^I \ar[r]_{p_0} \ar[d]^{p^I} & A \ar[d]^p \\
  & B^I \ar[r]_{p_0} & B
}
\end{equation}
there exists a homotopy $h^{'}$ keeping the diagram commutative.

\begin{lem}
\label{lem:hurewiczreplace}
\mbox{}
\begin{enumerate}
\item
Suppose that $i$ is a Hurewicz cofibration and that the diagram 
\[
\xymatrix{
  & A \ar@{>->}[dl]_i \ar[dr]^{g} & \\
  B \ar[rr]^{f} & & C
}
\]
is homotopy commutative $fi \simeq g$. Then there exists a map $f^{'} \simeq f$ with $f^{'}i = g$.
\item
Suppose that $p$ is a Hurewicz fibration and that the diagram
\[
\xymatrix{
  A \ar[rr]^{f} \ar[dr]_g & & B \ar@{->>}[dl]^{p} \\
  & C &
}
\]
is homotopy commutative $pf \simeq g$. Then there exists a map $f^{'} \simeq f$ with $pf^{'} = g$.
\end{enumerate}
\end{lem}

\begin{proof}
We only prove (1). Denote $h \colon fi \simeq g$ a homotopy. Using HEP for $i$, we extend $h$ to a homotopy $h^{'} \colon I \times B \ra C$ with $h^{'} \comp (I \times i) = h$ and $h^{'}i_0 = f$. Then $f^{'} = h^{'}i_1 \colon B \ra C$ has the desired properties.
\end{proof}

\index{Lemma!Dold lema}
\begin{lem}[A. Dold]
\label{lem:dold}
\mbox{}
\begin{enumerate}
\item
Suppose that $i$, $i^{'}$ are Hurewicz cofibrations and $fi = i^{'}$
\[
\xymatrix{
  & A \ar@{>->}[dl]_i \ar@{>->}[dr]^{i^{'}} & \\
  B \ar[rr]^{f} & & C
}
\]
If $f$ is a homotopy equivalence, then it is a homotopy equivalence {\it under} $A$.
\item
Suppose that $p$, $p^{'}$ are Hurewicz fibrations and $p^{'}f = p$
\[
\xymatrix{
  A \ar@{->>}[rr]^{p} \ar[dr]_g & & B \ar@{->>}[dl]^{p^{'}} \\
  & C &
}
\]
If $f$ is a homotopy equivalence, then it is a homotopy equivalence {\it over} $C$.
\end{enumerate}
\end{lem}

\begin{proof}
We only prove (1), and we do that in several steps.

Step 1. It suffices to show that $f$ admits a right homotopy inverse $f^{'}$ under $A$. For in that case, by the same reasoning $f^{'}$ has a right homotopy inverse $f^{''}$ under $A$, which has itself a right homotopy inverse $f^{'''}$ under $A$. We therefore have $f \overset{A}{\simeq} f^{''}$ and $f^{'} \overset{A}{\simeq} f^{'''}$, from which we see that $f^{'}$ is also a left homotopy inverse of $f$ under $A$.

Step 2. To show that $f$ admits a right homotopy inverse $f^{'}$ under $A$, it suffices to assume that $i = i^{'}$ and $f \simeq 1_B$. Indeed, suppose that $f^{'} \colon C \ra B$ is a homotopy inverse of $f$ (so $f^{'}f \simeq 1_{B}$). By \lemmaref{lem:hurewiczreplace}, we can find $f^{''} \simeq f^{'}$ with $f^{''}i^{'} = i$. Then $f^{''}f \simeq 1_B$, so by our assumption $f^{''}f \overset{A}{\simeq} 1_B$, and $f^{''}$ is a right homotopy inverse under $A$ to $f$.

Step 3. In the commutative diagram with $f \simeq 1_B$
\[
\xymatrix{
  & A \ar@{>->}[dl]_i \ar@{>->}[dr]^{i} & \\
  B \ar[rr]^{f} & & B
}
\]
we'd like to show that $f \overset{A}{\simeq} 1_B$. Denote $f \xymatrix{\ar[r]|h &} 1_{B}$ the homotopy $h \colon f \simeq 1_B$. Using HEP for $i$, the homotopy $i \xymatrix{\ar[rr]|{h\comp{(I \times i)}} &&} i$ extends to a homotopy $1_{B} \xymatrix{\ar[r]|{h^{'}} &} f^{'}$, for some map $f^{'} \colon B \ra B$. We will show that $f^{'}f \overset{A}{\simeq} 1_B$, which completes the proof of Step3, and with it the proof of our Lemma.

The composite $f^{'}f \xymatrix{&& \ar[ll]|{h^{'} \comp (I \times f)}} f \xymatrix{\ar[r]|{h} &} 1_{B}$ defines a homotopy $f^{'}f \simeq 1_{B}$. This homotopy restricts via $i$ to the homotopy $i \xymatrix{&& \ar[ll]|{h \comp (I \times i)}} i \xymatrix{\ar[rr]|{h \comp (I \times i)} &&} i$. The square
\[
\xymatrix{
  i \ar[d]|{id} & \; \ar@{}[drr]|{\alpha} & i \ar[ll]|{h \comp (I \times i)} \ar[rr]|{h \comp (I \times i)} && i \ar[d]|{id} \\
  i \ar[rrrr]|{id} &&& \; & i
}
\]
can be filled with a homotopy $\alpha \colon I^2 \times A \ra B$. We apply HEP to the Hurewicz cofibration $I \times A \ra I \times B$ with respect to the top edge of the previous diagram. This constructs a map $\beta \colon I^2 \times B \ra B$
\[
\xymatrix{
  f^{'}f \ar[d] & \; \ar@{}[drr]|{\beta} & f \ar[ll]|{h^{'} \comp (I \times f)} \ar[rr]|{h} && 1_{B} \ar[d] \\
  \, \ar[rrrr] &&& \; & \, 
}
\]
Tracing the left, bottom and right edges above yields the desired homotopy $f^{'}f \overset{A}{\simeq} 1_B$.
\end{proof}

\begin{lem}
\label{lem:trivhurcof}
\mbox{}
\begin{enumerate}
\item
Any trivial Hurewicz cofibration $i \colon A \ra B$ admits a strong deformation retract, i.e. a map $r \colon B \ra A$ such that $ri = 1_A$ and $ir \overset{A}{\simeq} 1_B$.
\item
Any trivial Hurewicz fibration $p \colon A \ra B$ admits a strong deformation section, i.e. a map $s \colon B \ra A$ such that $ps = 1_B$ and $sp \underset{A}{\simeq} 1_A$.
\end{enumerate}
\end{lem}

\begin{proof}
We only prove (1). Denote $r^{'}$ a homotopy inverse of $i$. By \lemmaref{lem:hurewiczreplace} applied to the homotopy commutative diagram
\[
\xymatrix{
  & A \ar@{>->}[dl]_i \ar[dr]^{1_A} & \\
  B \ar[rr]^{r^{'}} & & A
}
\]
there exists $r \colon B \ra A$, with $r \simeq r^{'}$ and $ri = 1_A$. From $ir \simeq 1_B$ in the commutative diagram
\[
\xymatrix{
  & A \ar@{>->}[dl]_i \ar@{>->}[dr]^{i} & \\
  B \ar[rr]^{ir} & & B
}
\]
by Dold's \lemmaref{lem:dold} we have $ir \overset{A}{\simeq} 1_B$.
\end{proof}

\begin{lem}
\label{lem:hurpush}
\mbox{}
\begin{enumerate}
\item
(Trivial) Hurewicz cofibrations are stable under pushouts.
\item
(Trivial) Hurewicz fibrations are stable under pullbacks.
\end{enumerate}
\end{lem}

\begin{proof}
We only prove (1). Hurewicz cofibrations are defined as having the left lifting property with respect to all $p_0 \colon B^I \ra B$, so Hurewicz cofibrations are stable under pushouts. Strong deformation retracts are also stable under pushouts, and in view of \lemmaref{lem:trivhurcof} so are trivial Hurewicz cofibrations.
\end{proof}

\begin{lem}
\label{lem:hurcolim}
\mbox{}
\begin{enumerate}
\item
If $A_0 \rightarrowtail A_1  \rightarrowtail ... \rightarrowtail A_n  \rightarrowtail ...$ is a sequence of (trivial) Hurewicz cofibrations, then its composition $A_0 \ra \colim^n A_n$ is again a (trivial) Hurewicz cofibration.
\item
If $... \ra A_n \ra ... \ra A_1 \ra A_0$ is a sequence of (trivial) Hurewicz fibrations, then $\lim^n A_n \ra A$ is again a (trivial) Hurewicz fibration.
\end{enumerate}
\end{lem}

\begin{proof}
We only prove (1). Denote $i_n \colon A_n \rightarrowtail A_{n+1}$, and $i \colon A_0 \ra \colim^n A_n = A$. If all $i_n$ are Hurewicz cofibrations, then $i$ has the left lifting property with respect to all $p_0 \colon B^I \ra B$, so $i$ is also a Hurewicz cofibration.

Suppose now that each $i_n$ is trivial. By \lemmaref{lem:trivhurcof}, there exist strong deformation retracts $r_n \colon A_{n+1} \ra A_n$ with $r_ni_n = 1_{A_n}$, and with relative homotopies $1_{A_{n+1}} \overset{A_n}{\simeq} i_nr_n$.

Define $r \colon A \ra A_0$ the colimit of the compositions $r_0r_1 ... r_n$. It is a retract of $i$, and we need to show that $1_A \simeq ir$.

The homotopy $1_{A_1} \overset{A_0}{\simeq} i_0r_0$ extends by HEP for $A_1 \rightarrowtail A$ to a homotopy $h_0 \colon 1_{A} \overset{A_0}{\simeq} s_0$, for some map $s_0 \colon A \ra A$ with $s_0 |_{A_1} = i_0r_0$.

The composite homotopy $s_0 |_{A_2} \simeq 1_{A_2} \simeq i_1i_0r_0r_1$, by Dold's \lemmaref{lem:dold} yields a homotopy $s_0 |_{A_2} \overset{A_1}{\simeq} i_1i_0r_0r_1$, which extends by HEP for $A_2 \rightarrowtail A$ to a homotopy $h_1 \colon s_0 \overset{A_1}{\simeq} s_1$, for some map $s_1 \colon A \ra A$ with $s_1 |_{A_2} = i_1i_0r_0r_1$.

By induction, we construct maps $s_n \colon A \ra A$ with $s_n |_{A_{n+1}} = i_n...i_0$$r_0...r_n$, and homotopies $h_n \colon s_{n-1} \overset{A_{n}}{\simeq} s_n$. Stitching together all the $h_n$, we obtain the desired homotopy $h \colon 1_A \simeq ir$ as
\[
h(1 - \frac{1}{2^n} + \frac{t}{2^{n+1}}, a) = h_n(t, a) \;\;\; (n \ge 0, \; 0 \le t \le 1, \; a \eps A)
\]
\end{proof}


\begin{thm}[The Hurewicz model structure]
\label{thm:hurtop}
$Top$ is an ABC model category, with all objects cofibrant and fibrant, with:
\begin{enumerate}
\item
Homotopy equivalences as weak equivalences
\item
Hurewicz fibrations as fibrations
\item
Hurewicz cofibrations as cofibrations.
\end{enumerate}
\end{thm}

\begin{proof}
The axioms CF1, CF2, CF5 are straightforward. The axiom CF4 is given by the classic mapping cylinder construction in $Top$. The axiom CF3 is proved by \lemmaref{lem:hurpush}, and CF6 by \lemmaref{lem:hurcolim}. All maps $p_0 \colon B^I \ra B$ admit the constant map $B \ra B^I$ as a section, from which any space is Hurewicz cofibrant. A similar proof shows that any space is also Hurewicz fibrant.
\end{proof}

We will mention without proof two more results regarding model structures on $Top$.

\begin{thm}[Strom, \cite{Strom1}]
\label{thm:stromtop}
$Top$ is a Quillen model category, with:
\begin{enumerate}
\item
Homotopy equivalences as weak equivalences
\item
Hurewicz fibrations as fibrations
\item
Hurewicz cofibrations that are closed maps, as cofibrations.
\end{enumerate}
\end{thm}

\index{Serre fibration}
\begin{defn}
\label{defn:serrefib}
\mbox{}
A map $p \colon A \ra B$ is a {\it Serre fibration}  if it has the right lifting property with respect to all face maps  $I^{n-1} \ra I^n$.
\end{defn}

\begin{thm}[Quillen, \cite{Quillen1}]
\label{thm:quillentop}
$Top$ is a Quillen model category, with:
\begin{enumerate}
\item
Weak homotopy equivalences as weak equivalences
\item
Serre fibrations as fibrations
\item
Maps with the left lifting property with respect to trivial Serre fibrations as cofibrations.
\end{enumerate}
\end{thm}

\section{Abelian categories}
\label{sec:exampleabcat}

We recall the following basic results regarding abelian categories:

\comment{
Suppose that ${\cat A}$ is an abelian category. This means that 

\begin{description}
\item[A0]
${\cat A}$ is pointed
\item[A1]
${\cat A}$ is closed under finite sums and products
\item[A3]
Every map has a kernel and a cokernel
\item[A4]
Every monic is a kernel of a map. Every epic is a cokernel of a map.
\end{description}

From these axioms, one shows \cite{Freyd1} that:
\begin{enumerate}
\item
${\cat A}$ is closed under finite limits and colimits
\item
For any pair of objects, the natural map $A \Sum B \ra A \times B$ is an isomorphism, whose common value is denoted $A \oplus B$
\item
${\cat A}$ is additive, i.e. ${\cat A}(A, B)$ is a natural abelian group with composition induced by the diagonal $B \oplus B \ra B$
\item
A map $f \colon A \ra B$ is monic iff $Ker \; f = \initial$, and is epic iff $Coker \; f = \initial$
\end{enumerate}
} %end comment


\begin{lem}
\label{lem:pushabelian1}
In an abelian category, given a square
\[
\xymatrix{
  A \ar[r]^f \ar[d]_g & B \ar[d]^{g^{'}} \\
  C \ar[r]^{f^{'}} & D
}
\]
consider the sequence $\xymatrix{A \ar[r]^-{(f, g)} & B \oplus C \ar[r]^-{f^{'} - g^{'}} & D}$. Then:
\begin{enumerate}
\item
$A \ra B \oplus C \ra D$ is zero iff the square commutes
\item
$A \ra B \oplus C \ra D \ra 0$ is exact iff the square is a pushout
\item
$0 \ra A \ra B \oplus C \ra D $ is exact iff the square is a pullback
\item
$0 \ra A \ra B \oplus C \ra D \ra 0$ is exact iff the square is a pushout and a pullback
\item
If the square is a pushout with $f$ or $g$ monic, then it is also a pullback
\item
If the square is a pullback with $f^{'}$ or $g^{'}$ epic, then it is also a pushout $\square$
\end{enumerate}
\end{lem}

\begin{proof}
Left to the reader (who may wish consult for example \cite{Freyd1}).
\end{proof}

\begin{lem}
\label{lem:pushabelian2}
In an abelian category:
\begin{enumerate}
\item
The pushout of a map $f$ is epic iff $f$ is epic.
\item
The pushout of a monic is monic.
\item
The pullback of a map $f^{'}$ is monic iff $f^{'}$ is monic.
\item
The pullback of an epic is epic $\square$
\end{enumerate}
\end{lem}


\begin{proof}
This is a ready consequence of \lemmaref{lem:pushabelian1}.
\end{proof}

In fact, any small abelian category admits, by a theorem of Mitchell, a fully faithful exact functor to a category of modules. As a consequence, to prove any statement involving small diagrams, (co)kernels, (co)images, monics, epics, pushouts, pullbacks and finite sums (like \lemmaref{lem:pushabelian1} and \lemmaref{lem:pushabelian2}) in an arbitrary abelian category ${\cat A}$ it suffices to prove that same statement for categories of modules.

\subsection*{Complexes of objects}
\index{$C({\cat A}), C^+({\cat A}), C^-({\cat A}), C^b({\cat A})$}
Denote $C({\cat A})$ the category of $\mathbb{Z}$-graded complexes of objects in ${\cat A}$ 
\begin{center}
$... \lra A_{n+1} \overset{a_{n+1}}{\lra} A_n \overset{a_{n}}{\lra} A_{n-1} \lra ...$

$a_n a_{n+1} = 0$ for any $n \eps {\mathbb Z}$
\end{center}
Denote $C^+({\cat A})$, $C^-({\cat A})$, $C^b({\cat A})$ the full subcategories of $C({\cat A})$ having as objects the bounded below, bounded above resp. bounded complexes. All four are abelian categories, with (co)kernels and (co)images computed degreewise. We recall the classic

\begin{lem}[Snake lemma]
\label{lem:snake}
\index{lemma!snake lemma}
Any short exact sequence $0 \ra A \sdot \overset{f}{\ra} B \sdot \overset{g}{\ra} C \sdot \ra 0$ in $C({\cat A})$ induces a natural long exact sequence in homology
\[
... \lra H_n A \overset{H_nf}{\lra} H_n B \overset{H_ng}{\lra} H_n C \overset{\delta_n}{\lra} H_{n-1} A \lra ...
\]
By Mitchell's Theorem, it suffices to define the connecting homomorphism $\delta_n$ for the case when ${\cat A}$ is a category of modules, in which case in the next diagram
\[
\xymatrix{
  0 \ar[r] & A_n \ar[d]_{a_n} \ar[r]^{f_n} & B_n \ar[d]^{b_n} \ar[r]^{g_n} & C_n \ar[d]^{c_n} \ar[r] & 0 \\
  0 \ar[r] & A_{n-1} \ar[r]^{f_{n-1}} & B_{n-1} \ar[r]^{g_{n-1}} & C_{n-1} \ar[r] & 0
}
\]
$\delta_n$ sends the representative of an element $x_n \eps Ker \, c_n$ to the representative of $y_{n-1} \eps Ker \, a_{n-1}$, where $y_{n-1}$  is constructed observing that there exists $z_n \eps B_n$ with $g_nz_n = x_n$, so $g_{n-1}b_n z_n = 0$ which allows us to define $y_{n-1} = f^{-1}_{n-1}b_nz_n$. Furthermore, the definition of $\delta_n$ outlined above does not depend on the choices involved. $\square$
\end{lem}

\index{quasi-isomorphism}
A map of complexes $A \sdot \ra B \sdot$ is called a {\it quasi-isomorphism} if it induces an isomorphism in homology. The homotopy category of $C({\cat A})$ (resp. $C^{+,-,b}({\cat A})$) with respect to quasi-isomorphisms is denoted $D({\cat A})$ (resp. $D^{+,-,b}({\cat A})$), and is called the {\it derived} category of ${\cat A}$.

\begin{thm}
\label{thm:abcatisabccat1}
For an abelian category ${\cat A}$, the categories of chain complexes $C({\cat A})$ (resp. $C^{+,-,b}({\cat A})$), with quasi-isomorphisms as weak equivalences, with monics as cofibrations and with epics as fibrations are all pointed ABC premodel categories, with all objects at once fibrant and cofibrant.
\end{thm}

\begin{proof}
Axiom CF3 (1) is given by \lemmaref{lem:pushabelian2} (2). Axiom CF3 (2) follows from the Snake lemma. To prove CF4, we construct the factorization of a map $f \colon A \sdot \ra B \sdot$ as a monic $f^{'}$ followed by a quasi-isomorphism $r$, as follows:
\[
\xymatrix {
  A \sdot \ar@{>->}[r]^-{f^{'}} & A \sdot \oplus A \sdot[1] \oplus B \sdot \ar[r]^-r_-\sim & B \sdot
}
\]
In the middle complex, the boundary map
\[
A_n \oplus A_{n-1} \oplus B_n \ra A_{n-1} \oplus A_{n-2} \oplus B_{n-1}
\]
is given by the matrix
\[
\begin{pmatrix}
  a_n & 0 & 0 \\
  id & -a_{n-1} & 0  \\
  0 & 0 & b_n 
\end{pmatrix}
\]

The ABC prefibration category axioms are proved in a dual fashion.
\end{proof}

\subsection*{Exact categories}
For an abelian category ${\cat A}$, suppose that ${\cat E} \subset {\cat A}$ is an exact subcategory, i.e. a full subcategory with the property that for any exact sequence
\[
0 \ra A \ra B \ra C \ra 0
\]
if two out of $A, B, C$ are in ${\cat E}$, then so is the third. Denote $C({\cat E})$ (resp. $C^{+, -, b}({\cat E})$) the category of complexes (resp. bounded below, bounded above, bounded complexes) of maps in ${\cat E}$. 

We say that a monic (resp. epic) map $f$ of ${\cat A}$ is {\it admissible}, if it is also a map of ${\cat E}$. We can now state the following stronger variant of \theoremref{thm:abcatisabccat1}:

\begin{thm}
\label{thm:abcatisabccat2}
Suppose that ${\cat A}$ is an abelian category and that ${\cat E} \subset {\cat A}$ is an exact subcategory. Then $C({\cat E})$ (resp. $C^{+,-,b}({\cat E})$), with quasi-isomorphisms as weak equivalences, with admissible monics as cofibrations and admissible epics as fibrations are all pointed ABC premodel categories, with all objects at once fibrant and cofibrant.
\end{thm}

\begin{proof}
Axiom CF3 (1) follows from \lemmaref{lem:pushabelian2} (2), observing that the push in ${\cat E}$ of an admissible monic exists and is again an admissible monic. Axioms CF3 (2) and CF4 are proved the same way as for \theoremref{thm:abcatisabccat1}.
\end{proof}

\subsection*{The Grothendieck axioms AB4, AB5} 
These are axioms that specify exactness properties for small sums, resp. for filtered colimits. Our reference here is Grothendieck's {\it Toh\^{o}ku} paper \cite{Grothendieck4}, also explained at length in \cite{Popescu1}. We start by recalling a classic adjunction lemma for abelian categories. 

\begin{lem}
\label{lem:adjabelian}
Suppose that $u \dashv v$, $u \colon {\cat A} \rightleftarrows {\cat B} \colon v$ is a pair of adjoint functors between two abelian categories. Then $u$ is right exact, and $v$ is left exact.
\end{lem}

\begin{proof}
Suppose that $A_0 \overset{f}{\ra} A_1 \overset{g}{\ra} A_2 \ra 0$ is an exact sequence in ${\cat A}$. Suppose that $B \eps {\cat B}$ is an object and $h \colon uA_1 \ra B$ is a map with $h \comp uf = 0$. 
\begin{equation}
\label{eqn:adjabelian1}
\xymatrix{
  uA_0 \ar[r]^{uf} & uA_1 \ar[r]^{ug} \ar[rd]_{h} & uA_2 \ar@{-->}[d]^{k^{'}} \ar[r] & 0 \\
  & & B &
}
\end{equation}
Denote $h^{'} \colon A_1 \ra vB$ the map adjoint to $h$. We have $h^{'}f = 0$, so there exists an unique map $k$ making the next diagram commutative,
\begin{equation}
\label{eqn:adjabelian2}
\xymatrix{
  A_0 \ar[r]^f & A_1 \ar[r]^g \ar[rd]_{h^{'}} & A_2 \ar@{-->}[d]^k \ar[r] & 0 \\
  & & vB &
}
\end{equation}
hence there exists a unique map $k^{'}$ (the adjoint of $k$) making the diagram \eqref{eqn:adjabelian1}  commutative. This shows that $ug = Coker(uf)$, so the top row of \eqref{eqn:adjabelian1} is exact.
\end{proof}

If ${\cat D}$ is a small category and ${\cat A}$ is abelian, then the category of ${\cat D}$-diagrams in ${\cat A}$, denoted ${\cat A}^{\cat D}$, is again abelian, with (co)kernels and (co)images computed pointwise. If ${\cat A}$ is closed under colimits indexed by ${\cat D}$, then the functor 
\[
\colim\nolimits^{\cat D} \colon {\cat A}^{\cat D} \ra {\cat A}
\]
is right exact by \lemmaref{lem:adjabelian}.

Grothendieck's axiom AB4 states that ${\cat A}$ is closed under small sums, and that for any set ${\cat D}$ the sum functor $\oplus^{\cat D}$ is left exact. 

An abelian category is always closed under pushouts; if it is closed under small sums then it is closed under all small colimits. Grothendieck's axiom AB5 states that ${\cat A}$ is closed under small sums (therefore under small colimits), and that for any {\it filtered} small category ${\cat D}$ the functor $\colim\nolimits^{\cat D}$ is left exact. Axiom AB5 implies axiom AB4.

One can formulate dual axioms AB4$^*$, AB5$^*$ for products and filtered limits, however, as explained in loc. cit., if an abelian category satisfies AB4 and AB4$^*$ then it is trivial in the sense that all its objects are isomorphic with $0$.

The next theorem is now immediate:

\begin{thm}
\label{thm:abcatisabccat3}
If an abelian category ${\cat A}$ satisfies AB4 (resp. AB5), then $C({\cat A})$, with quasi-isomorphisms as weak equivalences and monics as cofibrations satisfies CF5 (resp CF6). $\square$
\end{thm}

We will mention without proof the following result of Tibor Beke \cite{Beke1}:
\begin{thm}
\label{thm:abcatisabccat4}
If an abelian category ${\cat A}$ satisfies AB5 and has a generator (i.e. an object $A$ with the property that ${\cat A}(A, -)$ is faithful), then $C({\cat A})$, with quasi-isomorphisms as weak equivalences, with monics as cofibrations, and maps with the RLP with respect to monics as fibrations, forms a Quillen model category.
\end{thm}

